\scanpage{75}
identifiant les pts des centres des disques
qui se correspondent, et en ponctuant
par les centres des autres disques.
De plus, on obtient, \deleted{\unsure{un|une}} pour chaque
point singulier $x$ de $X'''$, correspondant
à une paire $\{x_1,x_2\}$ de points de
$X''$ (centres des deux disques associés),
grâce à l'antiisomorphisme de ces
disques pointés, un élément \inserted{non nul} \deleted{de} (i.e. une base)
de
\[
T_{X'',x_1}\underset{\mathbb C}{\otimes}T_{X'',x_2}
\]
i.e. \deleted{\unsure{de|un}} élément de
\[
T_{X'',x_1}^{*}\mathbin{\wedge}_{\mathbb C^*}T_{X'',x_2}^{*}.
\]

On trouve ainsi un foncteur de

a) Catégorie \inserted{(groupoïde)} des courbes hol. cptes à bord \inserted{$X$} de
type $g,\nu$, avec découpage holomorphe
(avec comme flèches les iso)

vers

b) Catégorie des courbes \inserted{$\mathbb C$-alg} stables de type $(g,\nu)$
$X'''$, munies \inserted{1°)} d'une ``structure de propagation
normale'' en ses pts singuliers, et 2°) d'un
\deleted{structure de multidisques} \inserted{syst.} admissible de
disques autour des pts marqués.
% archive sheet number: 77
